% =============================================================================
% Exercise list R3 -- worked solutions (MM benchmark, taxes, distress, trade-off)
% Statements: Listas/lista-r3-mm-taxes-tradeoff.md (not restated here).
% Numbers: Listas/lista-r3-mm-taxes-tradeoff-solutions_check.py (python -I, prints PASS).
% Compile: pdflatex (2x); check pdffonts (Type 1 only, uni = yes).
% =============================================================================
\documentclass[11pt,a4paper]{article}
\usepackage[utf8]{inputenc}
\usepackage[T1]{fontenc}
\usepackage{lmodern}   % Type 1 vector fonts (replaces PK bitmaps)
\usepackage{cmap}      % ToUnicode CMap -> copyable text and accents
\usepackage[brazil,english]{babel}
\usepackage{amsmath,amssymb,amsfonts}
\usepackage[a4paper,margin=2.2cm]{geometry}
\usepackage{booktabs}
\usepackage{enumitem}
\usepackage{fancyhdr}
\usepackage{xcolor}
\usepackage{microtype}

\pagestyle{fancy}\fancyhf{}
\lhead{\footnotesize Finanças II -- Corporativas (MPE Insper)}
\rhead{\footnotesize List R3 --- Worked solutions}
\cfoot{\footnotesize\thepage}
\renewcommand{\headrulewidth}{0.3pt}
\setlength{\headheight}{14pt}
\setlist{topsep=2pt,itemsep=1pt}

\makeatletter
\renewcommand\@seccntformat[1]{\csname the#1\endcsname.\quad}
\makeatother
\numberwithin{equation}{section}

\newcommand{\notation}{\par\smallskip\noindent\textbf{Notation.}}
\newcommand{\step}[2]{\par\smallskip\noindent\textbf{Step #1 --- #2.}~}
\newcommand{\intuition}{\par\smallskip\noindent\textbf{Intuition.}~}
\newcommand{\assumption}{\par\smallskip\noindent\textbf{Assumption.}~}
\newcommand{\bookans}{\par\smallskip\noindent\textbf{Book answer.}~}
\newcommand{\refbib}[1]{\par\smallskip\noindent{\footnotesize\color{black!65}\textbf{Ref:} #1}\par}
\newcommand{\E}{\mathbb{E}}

\begin{document}
\begin{center}
  {\Large\bfseries Finanças II -- Corporativas (MPE Insper)}\\[3pt]
  {\large Exercise list R3 --- A frictionless benchmark, taxes, distress and static trade-off}\\[3pt]
  \textit{Worked solutions}
\end{center}

\noindent\small Statements are in \texttt{Listas/lista-r3-mm-taxes-tradeoff.md} and are not
restated; each section cites the exercise by number and printed page. Editions cited are the
local copies: Berk, DeMarzo \& Harford, \textit{Fundamentals of Corporate Finance}, 2nd ed.\
(2011); Ross, Westerfield, Jordan \& Lamb, \textit{Fundamentos de Administração Financeira},
9th ed.\ (2013); Koski, Higgins \& Mitton, \textit{Analysis for Financial Management}, 13th
ed.\ (2023); Grinblatt \& Titman, \textit{Financial Markets and Corporate Strategy}, 2nd ed.\
(2002); Amaro de Matos, \textit{Theoretical Foundations of Corporate Finance} (2001). Every
number below is reproduced by \texttt{lista-r3-mm-taxes-tradeoff-solutions\_check.py}.
Amounts are in dollars unless a section says ``millions''.\normalsize

% =============================================================================
\section{Two-outcome project financed with equity or safe debt --- Berk 16.1, p.~490}

\notation
\begin{itemize}
  \item $C_{\text{low}}=130{,}000$, $C_{\text{high}}=180{,}000$: project free cash flow at $t=1$ in the two states; $p=\tfrac12$: probability of the low state.
  \item $I=100{,}000$: initial investment; $r_U=0.20$: project (unlevered) cost of capital; $r_f=0.10$: risk-free rate.
  \item $V_U$: value of the project's cash flows (= unlevered equity); $\mathit{NPV}=V_U-I$.
  \item $D=100{,}000$: amount borrowed at $r_f$; $E_L$: value of levered equity; $r_E$: its expected return; $r_D=r_f$.
\end{itemize}

\step{1}{Expected cash flow}
\begin{align}
  \E[C] &= p\,C_{\text{low}}+(1-p)\,C_{\text{high}} \label{eq:1EC}\\
        &= (0.5)(130{,}000)+(0.5)(180{,}000) \nonumber\\
        &= 155{,}000. \nonumber
\end{align}

\step{2}{(a) NPV} Discounting \eqref{eq:1EC} at the project's cost of capital,
\begin{align}
  V_U &= \frac{\E[C]}{1+r_U} \label{eq:1VU}\\
      &= \frac{(155{,}000)}{(1.20)} \nonumber\\
      &= 129{,}166.67, \nonumber\\
  \mathit{NPV} &= V_U - I \label{eq:1NPV}\\
      &= (129{,}166.67)-(100{,}000) \nonumber\\
      &= 29{,}166.67. \nonumber
\end{align}

\step{3}{(b) Unlevered equity} All-equity holders receive exactly $C_s$, so by the law of one
price their claim is worth \eqref{eq:1VU}: $E_U=V_U=129{,}166.67$. Selling the equity raises
\$129{,}166.67, more than $I$; the surplus is the NPV.

\step{4}{(c) Levered equity cash flows} The safe loan requires a repayment of
\begin{align}
  D(1+r_f) &= (100{,}000)(1.10) \nonumber\\
           &= 110{,}000, \label{eq:1rep}
\end{align}
which is below the lowest $C_s$, so the debt is indeed riskless. Equity receives
$C-110{,}000$, i.e.\ $130{,}000-110{,}000=20{,}000$ or $180{,}000-110{,}000=70{,}000$.

\step{5}{(c) Value by MM I} With no taxes or distress costs, MM I gives $E_L+D=V_U$, so
\begin{align}
  E_L &= V_U - D \label{eq:1EL}\\
      &= (129{,}166.67)-(100{,}000) \nonumber\\
      &= 29{,}166.67. \nonumber
\end{align}

\step{6}{Consistency with MM II} Expected equity cash flow is $(0.5)(20{,}000)+(0.5)(70{,}000)=45{,}000$, so
\begin{align}
  r_E &= \frac{45{,}000}{E_L}-1 \label{eq:1rE}\\
      &= \frac{45{,}000}{(29{,}166.67)}-1 \nonumber\\
      &= 0.542857, \nonumber
\end{align}
and MM II, $r_E=r_U+\tfrac{D}{E_L}(r_U-r_D)$, with $\tfrac{D}{E_L}=\tfrac{100{,}000}{29{,}166.67}=3.428571$, gives
\begin{align}
  r_E &= (0.20)+(3.428571)(0.20-0.10) \nonumber\\
      &= 0.542857, \nonumber
\end{align}
the same as \eqref{eq:1rE}: levered equity is worth \$29{,}166.67 and earns 54.29\%.

\intuition Borrowing slices the same cash flows differently; it does not change their total
value. The entrepreneur captures the same NPV, \$29{,}166.67, either as the excess of the
equity sale over $I$ or as the value of the levered equity left after the loan. The riskier
levered equity simply demands a higher expected return.
\refbib{Berk Ch.~16, §16.2 ``Capital Structure in Perfect Capital Markets'', pp.~464--472 (MM I, MM II).}

% =============================================================================
\section{One-time dividend or repurchase in a perfect market --- Berk 17.7, p.~528}

\notation
\begin{itemize}
  \item $N=500$ million shares; $P_0=15$: cum-payout share price; $C=250$ million: excess cash paid out.
  \item $d$: dividend per share; $P_{\text{ex}}$: ex-dividend price; $n_R$: shares repurchased (millions); $P_R$: post-repurchase price.
\end{itemize}

\step{1}{Market value of the firm} Equity (the firm has no debt) is worth
$N P_0=(500)(15)=7{,}500$ million, cash included.

\step{2}{(a) Ex-dividend price}
\begin{align}
  d &= \frac{C}{N}=\frac{(250)}{(500)}=0.50, \label{eq:2d}\\
  P_{\text{ex}} &= \frac{N P_0 - C}{N} \label{eq:2Pex}\\
     &= \frac{(7{,}500)-(250)}{(500)} \nonumber\\
     &= 14.50. \nonumber
\end{align}

\step{3}{(b) Repurchase} In a perfect market the shares are bought at the fair price $P_0$:
\begin{align}
  n_R &= \frac{C}{P_0}=\frac{(250)}{(15)}=16.6667, \label{eq:2nR}\\
  P_R &= \frac{N P_0 - C}{N-n_R} \label{eq:2PR}\\
      &= \frac{(7{,}500)-(250)}{(500)-(16.6667)} \nonumber\\
      &= \frac{7{,}250}{483.3333} \nonumber\\
      &= 15.00. \nonumber
\end{align}

\step{4}{(c) Comparison} Per original share, the dividend leaves $P_{\text{ex}}+d=(14.50)+(0.50)=15.00$
in stock plus cash, by \eqref{eq:2d}--\eqref{eq:2Pex}; the repurchase leaves either a share worth
15.00 by \eqref{eq:2PR} or 15.00 in cash for a tendered share. Wealth is \$15 per share under
both policies: neither makes investors better off.

\intuition Paying out cash shrinks the firm by exactly the cash paid. With a dividend the price
falls by the dividend; with a repurchase the share count falls instead and the price stays put.
Any investor can create the other pattern for herself (sell shares, or reinvest the dividend),
so payout policy is irrelevant (MM).
\refbib{Berk Ch.~17, §17.2 ``Dividends Versus Share Repurchases in a Perfect Capital Market'', pp.~503--508.}

% =============================================================================
\section{MM I with corporate taxes --- Ross 16.14, p.~578}

\notation
\begin{itemize}
  \item $\mathit{EBIT}=92{,}000$ per year forever; $R_U=0.15$: unlevered cost of equity; $R_D=0.09$: borrowing rate.
  \item $T_C=0.35$: corporate tax rate; $D=60{,}000$: perpetual debt used to buy back shares.
  \item $V_U$, $V_L$: unlevered and levered firm values; $E$: levered equity value.
\end{itemize}

\step{1}{Unlevered value} Cash flow to the all-equity firm is $\mathit{EBIT}(1-T_C)=(92{,}000)(1-0.35)=59{,}800$ forever, so
\begin{align}
  V_U &= \frac{\mathit{EBIT}(1-T_C)}{R_U} \label{eq:3VU}\\
      &= \frac{(59{,}800)}{(0.15)} \nonumber\\
      &= 398{,}666.67. \nonumber
\end{align}

\step{2}{Value of the tax shield} The yearly saving is $T_C R_D D=(0.35)(0.09)(60{,}000)=1{,}890$, a
perpetuity as risky as the debt, hence discounted at $R_D$:
\begin{align}
  \mathit{PV}(\text{TS}) &= \frac{T_C R_D D}{R_D} \label{eq:3TS}\\
      &= T_C D \nonumber\\
      &= \frac{(1{,}890)}{(0.09)} \nonumber\\
      &= 21{,}000. \nonumber
\end{align}

\step{3}{Levered value (MM I with taxes)} Adding \eqref{eq:3TS} to \eqref{eq:3VU},
\begin{align}
  V_L &= V_U + T_C D \label{eq:3VL}\\
      &= (398{,}666.67)+(21{,}000) \nonumber\\
      &= 419{,}666.67, \nonumber
\end{align}
of which equity is $E=V_L-D=(419{,}666.67)-(60{,}000)=359{,}666.67$.

\intuition The government's claim on EBIT falls by \$1{,}890 a year once interest is
deductible; that perpetual saving, worth $T_C D$, accrues to the firm's investors. The
borrowing rate does not affect the answer because it scales both the saving and its discount
rate.
\refbib{Ross Cap.~16, §16.4 ``Proposições I e II de M\&M com impostos sobre o lucro da pessoa jurídica'', pp.~551--556.}

% =============================================================================
\section{APV with constant debt --- Higgins 8.23, p.~333}

\notation
\begin{itemize}
  \item $\mathit{NPV}_0=500{,}000$: base-case NPV of the project's free cash flows.
  \item $D=1{,}000{,}000$: constant debt for $T=3$ years; $r=0.08$: interest and discount rate of the shields; $t=0.40$: tax rate.
  \item $\mathit{TS}$: annual interest tax shield, received at the end of years 1--3; $\mathit{APV}$: adjusted present value.
\end{itemize}

\step{1}{Annual shield}
\begin{align}
  \mathit{TS} &= D\,r\,t \label{eq:4TS}\\
       &= (1{,}000{,}000)(0.08)(0.40) \nonumber\\
       &= 32{,}000. \nonumber
\end{align}

\step{2}{Three-year annuity factor at 8\%}
\begin{align}
  A &= \frac{1}{1.08}+\frac{1}{1.08^2}+\frac{1}{1.08^3} \label{eq:4A}\\
    &= (0.925926)+(0.857339)+(0.793832) \nonumber\\
    &= 2.577097. \nonumber
\end{align}

\step{3}{PV of the shields and APV} Using \eqref{eq:4TS} and \eqref{eq:4A},
\begin{align}
  \mathit{PV}(\text{TS}) &= \mathit{TS}\cdot A \label{eq:4PV}\\
      &= (32{,}000)(2.577097) \nonumber\\
      &= 82{,}467.10, \nonumber\\
  \mathit{APV} &= \mathit{NPV}_0+\mathit{PV}(\text{TS}) \label{eq:4APV}\\
     &= (500{,}000)+(82{,}467.10) \nonumber\\
     &= 582{,}467.10. \nonumber
\end{align}

\bookans Matches the book's answer, \$582{,}467 (p.~425).

\intuition APV separates the value of the operations (\$500{,}000) from the value created by the
financing side effect. With a fixed debt the shields are as safe as the debt, so they are
discounted at the debt rate.
\refbib{Higgins Ch.~8, Appendix ``Asset Beta and Adjusted Present Value'', pp.~318--325 (APV, p.~321).}

% =============================================================================
\section{MM II with a rising cost of debt --- Berk 16.5, p.~491}

\notation
\begin{itemize}
  \item $r_U=0.12$: expected return on unlevered equity (= on the assets, unchanged by leverage without taxes).
  \item $D/E$: debt-equity ratio (0.50 or 1.50); $r_D$: debt cost of capital (0.06 or 0.08); $r_E$: expected return on levered equity.
\end{itemize}

\step{1}{MM II} With no taxes, the asset return is the value-weighted average of the claims,
$r_U=\tfrac{E}{D+E}r_E+\tfrac{D}{D+E}r_D$; solving for $r_E$,
\begin{align}
  r_E = r_U + \frac{D}{E}\,(r_U-r_D). \label{eq:5MM2}
\end{align}

\step{2}{(a) $D/E=0.50$, $r_D=6\%$} Substituting into \eqref{eq:5MM2},
\begin{align}
  r_E &= (0.12)+(0.50)(0.12-0.06) \nonumber\\
      &= 0.15. \label{eq:5a}
\end{align}

\step{3}{(b) $D/E=1.50$, $r_D=8\%$} Substituting into \eqref{eq:5MM2},
\begin{align}
  r_E &= (0.12)+(1.50)(0.12-0.08) \nonumber\\
      &= 0.18. \label{eq:5b}
\end{align}

\step{4}{Check: the WACC is unchanged} With $E/V=1/1.5$ and $D/V=0.5/1.5$ in (a), and $E/V=1/2.5$, $D/V=1.5/2.5$ in (b),
\begin{align}
  \text{(a)}\quad & \tfrac{1}{1.5}(0.15)+\tfrac{0.5}{1.5}(0.06)=0.12, \nonumber\\
  \text{(b)}\quad & \tfrac{1}{2.5}(0.18)+\tfrac{1.5}{2.5}(0.08)=0.12. \nonumber
\end{align}

\step{5}{(c) Reply to the manager} The higher expected return in \eqref{eq:5b} is pure
compensation for higher risk: equity bears the same business risk on a thinner base. By MM I the
firm's value, and therefore the share price, is the same at every debt level; the WACC stays at
12\%. Choosing leverage to maximise $r_E$ does not raise shareholder wealth; shareholders can
also lever up on their own (homemade leverage) if they want a riskier, higher-return position.
Shareholders care about the value of their claim, not its expected return.

\intuition The cost of equity rises linearly with $D/E$ at slope $r_U-r_D$. When debt becomes
riskier ($r_D$ rises), part of the risk moves to creditors and $r_E$ rises less steeply, but the
total required return on the assets is still $r_U$.
\refbib{Berk Ch.~16, §16.2, pp.~464--472 (MM II, leverage and the cost of equity).}

% =============================================================================
\section{Tax shields on amortising debt --- Berk 16.11, p.~492}

\notation
\begin{itemize}
  \item $B_t$: debt outstanding during year $t=1,\dots,4$ (millions); $B_1=100$, $B_{t+1}=B_t-25$.
  \item $r_D=0.10$: interest rate and discount rate of the shields; $\tau_c=0.40$: marginal corporate tax rate.
  \item $\mathit{TS}_t=\tau_c r_D B_t$: tax shield received at the end of year $t$.
\end{itemize}

\assumption The first \$25 million repayment is made one year from now, together with the
first interest payment, and interest in year $t$ is charged on the balance at the start of the
year. The loan is then repaid after four years.

\step{1}{Schedule} Balances are $100, 75, 50, 25$, so interest $r_D B_t$ is $10, 7.5, 5, 2.5$ and
\begin{align}
  \mathit{TS}_t &= (0.40)(0.10)B_t \in \{4,\ 3,\ 2,\ 1\}. \label{eq:6TS}
\end{align}

\step{2}{Present value} Discounting \eqref{eq:6TS} at $r_D$ (same risk as the loan),
\begin{align}
  \mathit{PV}(\text{TS}) &= \sum_{t=1}^{4}\frac{\mathit{TS}_t}{(1+r_D)^t} \label{eq:6PV}\\
      &= \frac{(4)}{1.10}+\frac{(3)}{1.10^2}+\frac{(2)}{1.10^3}+\frac{(1)}{1.10^4} \nonumber\\
      &= 3.636364+2.479339+1.502630+0.683013 \nonumber\\
      &= 8.301346. \nonumber
\end{align}
The interest tax shields are worth about \$8.30 million.

\intuition Each year's shield is proportional to the debt still outstanding, so amortisation
shrinks the shield stream. A permanent \$100 million loan would be worth $\tau_c D=\$40$
million in shields; repaying it over four years leaves only about a fifth of that.
\refbib{Berk Ch.~16, §16.3 ``Debt and Taxes'', pp.~473--479 (value of the interest tax shield).}

% =============================================================================
\section{Homemade leverage between twin firms --- Ross 16.18, p.~578}

\notation
\begin{itemize}
  \item $X=400{,}000$: perpetual operating income of each firm (random, identical for both); no taxes.
  \item Weber (unlevered): $E_W=V_W=2{,}500{,}000$. Cavaleiro (levered): $E_C=1{,}632{,}000$, $D_C=1{,}200{,}000$, interest $I_C=72{,}000$ per year, so $r_D=72{,}000/1{,}200{,}000=0.06$.
  \item $\alpha=0.05$: fraction bought; $r=0.06$: personal borrowing rate; $L$: amount borrowed personally.
\end{itemize}

\assumption ``Os custos líquidos das estratégias sejam iguais'' is read as: both strategies
require the same own funds today, the \$81{,}600 that 5\% of Cavaleiro costs.

\step{1}{Strategy C: 5\% of Cavaleiro} Equity income is $X-I_C=400{,}000-72{,}000=328{,}000$, so
\begin{align}
  \text{cost}_C &= \alpha E_C=(0.05)(1{,}632{,}000)=81{,}600, \label{eq:7cC}\\
  \text{income}_C &= \alpha (X-I_C)=(0.05)(328{,}000)=16{,}400, \label{eq:7iC}
\end{align}
a return of $16{,}400/81{,}600=20.10\%$.

\step{2}{Strategy W: 5\% of Weber with the same own funds} The stake costs $(0.05)(2{,}500{,}000)=125{,}000$
and earns $(0.05)(400{,}000)=20{,}000$. With own funds \eqref{eq:7cC},
\begin{align}
  L &= 125{,}000-81{,}600 = 43{,}400, \label{eq:7L}\\
  \text{income}_W &= 20{,}000 - rL = (20{,}000)-(0.06)(43{,}400) \nonumber\\
      &= (20{,}000)-(2{,}604) = 17{,}396, \label{eq:7iW}
\end{align}
a return of $17{,}396/81{,}600=21.32\%$.

\step{3}{(a) The return is higher at no extra risk} Both incomes contain $0.05X$, so they move
one-for-one with the same operating risk; \eqref{eq:7iW} exceeds \eqref{eq:7iC} by
$17{,}396-16{,}400=996$ in every state. Yes: the investor raises her return by buying Weber and
borrowing, and the position is no riskier (it even uses less leverage: \$43{,}400 of personal debt
against her 5\% share of Cavaleiro's \$1{,}200{,}000, i.e.\ \$60{,}000).

\step{4}{The underlying arbitrage} Replicating strategy C exactly needs $L=\alpha D_C=(0.05)(1{,}200{,}000)=60{,}000$:
income $20{,}000-(0.06)(60{,}000)=16{,}400$, identical to \eqref{eq:7iC}, for own funds
$125{,}000-60{,}000=65{,}000$. The saving is
\begin{align}
  81{,}600-65{,}000 = 16{,}600 = \alpha (V_C-V_W) = (0.05)(2{,}832{,}000-2{,}500{,}000), \label{eq:7arb}
\end{align}
and investing it at 6\% yields $(0.06)(16{,}600)=996$ a year, the gap found in Step~3.

\step{5}{(b) What investors do and when it stops} Every Cavaleiro shareholder gains by selling
Cavaleiro shares and buying Weber shares with personal borrowing. Selling pressure lowers $E_C$
and buying pressure raises $E_W$ until the arbitrage profit \eqref{eq:7arb} is zero, i.e.\ until
$V_C=V_W$ (MM I). If Weber's value stays at \$2.5 million, it stops when
$E_C=2{,}500{,}000-1{,}200{,}000=1{,}300{,}000$; Cavaleiro's cost of equity is then
$328{,}000/1{,}300{,}000=25.23\%$, which is MM II: $0.16+\tfrac{1.2}{1.3}(0.16-0.06)=25.23\%$,
with $r_U=400{,}000/2{,}500{,}000=0.16$.

\intuition Investors can borrow on the same terms as the firm, so corporate leverage adds
nothing they cannot build themselves. A levered firm that trades at a premium is selling
something investors can make more cheaply at home, and arbitrage removes the premium.
\refbib{Ross Cap.~16, §16.2 ``O efeito da alavancagem financeira'' (homemade leverage), pp.~542--547; §16.3, pp.~547--551 (Propositions I and II).}

% =============================================================================
\section{Announcement effect of a debt-financed repurchase with distress costs --- Berk 16.22, p.~494}

\notation
\begin{itemize}
  \item $P_0=5.50$, $N=10$ million shares, so $E_0=(5.50)(10)=55$ million.
  \item $D=20$ million of permanent new debt; $\tau_c=0.30$; $P_1$: price after the announcement; $\mathit{PV}(\text{FDC})$: PV of financial distress costs.
\end{itemize}

\assumption The announcement is credible and all the value change is priced at once; shares are
later bought back at the post-announcement price, so the repurchase itself transfers no value.

\step{1}{(a) Taxes only} Permanent debt adds $\tau_c D=(0.30)(20)=6$ million (cf.\ \eqref{eq:3TS}), all of it to the current shareholders:
\begin{align}
  P_1 &= \frac{E_0+\tau_c D}{N} \label{eq:8P1}\\
      &= \frac{(55)+(6)}{(10)} \nonumber\\
      &= 6.10. \nonumber
\end{align}

\step{2}{(b) Taxes and distress costs} By the trade-off identity $V_L=V_U+\tau_c D-\mathit{PV}(\text{FDC})$,
the observed equity gain is
\begin{align}
  (5.75)(10)-(55) &= 2.5 = \tau_c D-\mathit{PV}(\text{FDC}), \nonumber\\
  \mathit{PV}(\text{FDC}) &= (6)-(2.5) \nonumber\\
      &= 3.5 \text{ million}. \label{eq:8FDC}
\end{align}

\intuition The market pays the shareholders for the expected tax savings and charges them for
expected distress costs. The smaller price jump (\$0.25 against \$0.60) shows that the market
expects distress costs to eat more than half of the tax benefit of this debt.
\refbib{Berk Ch.~16, §16.3 pp.~473--479 and §16.5 ``Optimal Capital Structure: The Tradeoff Theory'', pp.~480--482.}

% =============================================================================
\section{After-tax WACC with riskless perpetual debt, then a debt-for-equity exchange --- Grinblatt \& Titman 13.15--13.16, p.~494}

\notation
\begin{itemize}
  \item $D=50$, $E=50$ (millions): perpetual riskless debt and equity; $A=D+E=100$: market value of the assets.
  \item $\beta_E=1.2$; $r_f=0.08$; $\E[r_M]=0.13$; CAPM holds; $r_D=r_f$ (riskless debt); $T_c=0.40$.
  \item $\mathit{UA}$, $\mathit{TX}$: unlevered assets and debt tax shield, $A=\mathit{UA}+\mathit{TX}$; $\beta_{UA}$, $r_{UA}$: their beta and expected return.
  \item Primes ($D'$, $E'$, $V'$, $\mathit{WACC}'$) denote values after the exchange offer.
\end{itemize}

\step{1}{13.15: cost of equity (CAPM)}
\begin{align}
  r_E &= r_f+\beta_E(\E[r_M]-r_f) \nonumber\\
      &= (0.08)+(1.2)(0.13-0.08) \nonumber\\
      &= 0.14. \label{eq:9rE}
\end{align}

\step{2}{13.15: WACC} G\&T's after-tax WACC (13.8) with \eqref{eq:9rE}:
\begin{align}
  \mathit{WACC} &= \frac{E}{D+E}\,r_E+\frac{D}{D+E}\,r_D(1-T_c) \label{eq:9W}\\
      &= (0.5)(0.14)+(0.5)(0.08)(1-0.40) \nonumber\\
      &= 0.094. \nonumber
\end{align}

\step{3}{Unlevered cost of capital (Hamada)} With static perpetual riskless debt, $\mathit{TX}=T_cD$ (13.2) and
$\mathit{UA}=D+E-T_cD=100-(0.40)(50)=80$ (13.3). Unlevering the equity beta with (13.7),
\begin{align}
  \beta_{UA} &= \frac{\beta_E}{1+(1-T_c)D/E} \label{eq:9bUA}\\
      &= \frac{(1.2)}{1+(1-0.40)(1)} \nonumber\\
      &= 0.75, \nonumber\\
  r_{UA} &= r_f+\beta_{UA}(\E[r_M]-r_f) \label{eq:9rUA}\\
      &= (0.08)+(0.75)(0.05) \nonumber\\
      &= 0.1175. \nonumber
\end{align}
Check with the adjusted cost of capital formula (p.~482), $\mathit{WACC}=r_{UA}\bigl(1-T_c\tfrac{D}{D+E}\bigr)$:
$(0.1175)(1-(0.40)(0.5))=0.094$, equal to \eqref{eq:9W}.

\step{4}{13.16: values after retiring \$25 million of debt} The unlevered assets do not change;
the shield shrinks to $T_cD'=(0.40)(25)=10$. Hence
\begin{align}
  V' &= \mathit{UA}+T_cD' = (80)+(10) = 90, \label{eq:9V}\\
  E' &= V'-D' = (90)-(25) = 65, \nonumber
\end{align}
of which \$25 million is the new shares given to the old creditors and $50-(0.40)(25)=40$ is the
old shareholders' stake: they bear the \$10 million loss of shield value.

\step{5}{13.16: new WACC} Substituting $D'/V'=25/90=0.277778$ into the adjusted cost of capital formula with \eqref{eq:9rUA},
\begin{align}
  \mathit{WACC}' &= r_{UA}\Bigl(1-T_c\frac{D'}{V'}\Bigr) \label{eq:9W2}\\
      &= (0.1175)\bigl(1-(0.40)(0.277778)\bigr) \nonumber\\
      &= (0.1175)(0.888889) \nonumber\\
      &= 0.104444. \nonumber
\end{align}

\step{6}{Cross-check through the cost of equity} By (13.6), $\beta_E'=\beta_{UA}\bigl(1+(1-T_c)\tfrac{D'}{E'}\bigr)
=(0.75)\bigl(1+(0.6)\tfrac{25}{65}\bigr)=(0.75)(1.230769)=0.923077$, so $r_E'=(0.08)+(0.923077)(0.05)=0.126154$ and
\begin{align}
  \mathit{WACC}' &= \tfrac{65}{90}(0.126154)+\tfrac{25}{90}(0.08)(0.6) \nonumber\\
      &= (0.091111)+(0.013333) = 0.104444, \nonumber
\end{align}
the same as \eqref{eq:9W2}. The WACC rises from 9.40\% to 10.44\%.

\assumption The debt ratio in \eqref{eq:9W2} is measured at post-exchange market values. Using
$D'/V=25/100$ (ignoring the lost shield) would give $(0.1175)(1-0.40\cdot0.25)=10.575\%$; that
mixes the old firm value with the new debt and is not consistent with the Hamada model.

\intuition In the Hamada model the only thing debt adds is a riskless tax shield. Cutting the debt
removes part of that cheap, government-funded slice of the financing, so the WACC moves back
toward the unlevered cost of capital, 11.75\%.
\refbib{Grinblatt \& Titman Ch.~13: Hamada model, eqs.\ (13.1)--(13.7), pp.~463--466; WACC (13.8), p.~476; adjusted cost of capital formula, pp.~481--482.}

% =============================================================================
\section{Taxes vs.\ bankruptcy costs in two states --- Grinblatt \& Titman 16.2, p.~591}

\notation
\begin{itemize}
  \item $C_G=210$, $C_B=66$: pre-tax cash flows next period in the favourable and unfavourable state, each with probability $\tfrac12$.
  \item $t=\tfrac13$: tax rate of the all-equity firm; $s=0.05$: tax saved per dollar of debt financing; $r_f=0.10$; investors are risk neutral.
  \item $F$: face value of the zero-coupon debt; $D$: its market value today; $K=20$: bankruptcy cost; $V_U$, $E$, $V$: unlevered value, equity, levered firm value.
\end{itemize}

\assumption ``Debt financing used'' is the amount raised today, $D$; the saving $sD$ is
realised next period in both states (if the firm defaults it accrues to the creditors who take
it over). The answer to (e) does not depend on this reading (Step~6).

\step{1}{(a) All-equity value} After-tax flows are $210(1-\tfrac13)=140$ and $66(1-\tfrac13)=44$, so
\begin{align}
  V_U &= \frac{\tfrac12(140)+\tfrac12(44)}{1.10} \label{eq:10VU}\\
      &= \frac{920}{11} \nonumber\\
      &= 83.6364. \nonumber
\end{align}

\step{2}{(b) Debt with $F=44$} Guess the debt is safe: $D=44/1.10=40$, so the saving is $sD=(0.05)(40)=2$.
In the bad state the firm then has $44+2=46\ge 44$: no default, confirming the guess. The firm
receives $D=40$.

\step{3}{(c) Equity} Equity gets $140+2-44=98$ (good) and $46-44=2$ (bad):
\begin{align}
  E &= \frac{\tfrac12(98)+\tfrac12(2)}{1.10} \label{eq:10E}\\
    &= \frac{500}{11} \nonumber\\
    &= 45.4545. \nonumber
\end{align}

\step{4}{(d) Firm value}
\begin{align}
  V &= D+E \label{eq:10V}\\
    &= (40)+(45.4545) \nonumber\\
    &= 85.4545, \nonumber
\end{align}
equal to $V_U+\tfrac{2}{1.10}=\tfrac{920}{11}+\tfrac{20}{11}=\tfrac{940}{11}$: the PV of the tax saving, $1.8182$, is the whole gain.

\step{5}{(e) Debt with $F=70$} Now $44+sD<70$ in the bad state, so the firm defaults; creditors
take over and receive $44+0.05D-20=24+0.05D$. In the good state they get 70. Hence
\begin{align}
  D &= \frac{\tfrac12(70)+\tfrac12(24+0.05D)}{1.10} \label{eq:10De}\\
  1.10\,D &= 47+0.025\,D \nonumber\\
  D &= \frac{47}{1.075} = 43.72093. \nonumber
\end{align}
With $sD=2.18605$, the bad-state recovery is $26.18605<70$ (default confirmed) and equity gets
$70+2.18605=72.18605$ in the good state and 0 in the bad state:
\begin{align}
  E &= \frac{\tfrac12(72.18605)}{1.10} \nonumber\\
    &= 32.81184, \nonumber\\
  V &= D+E \label{eq:10Ve}\\
    &= (43.72093)+(32.81184) \nonumber\\
    &= 76.53277. \nonumber
\end{align}
Equivalently $V=V_U+\mathit{PV}(sD)-\mathit{PV}(K)=(83.63636)+(1.98732)-(9.09091)=76.53277$, using $2.18605/1.10=1.98732$ and $\tfrac12(20)/1.10=9.09091$.
No: the firm is worth less with a \$70 obligation (76.53 against 85.45).

\step{6}{Robustness} Whatever the basis of the saving, it is at most $0.05\times70$ next period,
worth $3.18182$, while expected bankruptcy costs are worth $9.09091$; so
$V\le(83.63636)+(3.18182)-(9.09091)=77.72727<85.4545$ under any reading.

\intuition This is the static trade-off in miniature. Safe debt adds a small tax subsidy at no
cost. Debt large enough to default in the bad state buys a little more subsidy but triggers a
deadweight loss of \$20 with probability one half, which dominates.
\refbib{Grinblatt \& Titman Ch.~16, §16.1 ``Bankruptcy'' (direct bankruptcy costs vs.\ the debt tax advantage), pp.~559--561.}

% =============================================================================
\section{Risky debt in a binomial market and new debt of equal seniority --- Grinblatt \& Titman 14.2--14.3, pp.~526--527}

\notation
\begin{itemize}
  \item $C_G=210$, $C_B=80$: firm cash flow in the good/bad state; risk-neutral probability $\tfrac12$ each; $r_f=0$; no taxes.
  \item $F_0=50$: face of the existing debt; $F_1=50$: face of the new equal-seniority debt; 100 shares initially.
  \item $V$: firm (machine) value; $D_0$, $D_1$: market values of old and new debt; $P$: share price; $n_R$: shares repurchased.
\end{itemize}

\step{1}{14.2: share price} With $r_f=0$, values are risk-neutral expectations:
\begin{align}
  V &= \tfrac12(210)+\tfrac12(80) = 145, \label{eq:11V}
\end{align}
the old debt is safe ($80\ge50$), $D_0=50$, equity is $145-50=95$, and $P=95/100=0.95$.

\step{2}{(a) Price of the new debt} Total face becomes $F_0+F_1=100>80$, so in the bad state the
\$80 is shared pro rata, 40 to each issue. The new debt is worth
\begin{align}
  D_1 &= \tfrac12(50)+\tfrac12(40) = 45. \label{eq:11D1}
\end{align}
The old debt now has the same payoff, so $D_0$ also falls from 50 to 45.

\step{3}{(b) Share price at the announcement} The firm holds the machine plus the \$45 cash raised, and owes debt worth $2\times45=90$:
\begin{align}
  E &= (145)+(45)-(90) = 100, \qquad P = \frac{100}{100} = 1.00. \label{eq:11P}
\end{align}
The firm then buys back $n_R=45/1.00=45$ shares; 55 remain, and equity is $145-90=55$, paying 110
or 0, worth $\tfrac12(110)=55$, i.e.\ still \$1.00 per share.

\assumption ``A shareholder with 20 percent of the shares'' holds 20 of the 100 shares before
the announcement; ``undoing the leverage change'' means rebuilding the pre-announcement payoff
exposure with personal trades.

\step{4}{(c) Undoing the leverage} Before: 20\% of an equity paying $(160,30)$, i.e.\ $(32,6)$,
worth $\tfrac12(32)+\tfrac12(6)=19$. After the announcement the 20 shares are worth 20. Keep
$0.2\times55=11$ shares (paying $(22,0)$) and sell $20-11=9$ shares at \$1.00; with the \$9 buy
20\% of the new debt, $0.2\times45=9$, which pays $0.2\times(50,40)=(10,8)$. The position pays
\begin{align}
  (22+10,\ 0+8) = (32,\ 8), \qquad \text{worth } \tfrac12(32)+\tfrac12(8) = 20. \label{eq:11c}
\end{align}
This beats the old payoff $(32,6)$ state by state; the gain $20-19=1$ is 20\% of the \$5 lost by
the old creditors. Alternatively, she can rebuild $(32,6)$ exactly with $12.25$ shares and 7.5\%
of all the debt (paying $0.075\times(100,80)=(7.5,6)$), at a cost of $12.25(1.00)+0.075(90)=19.00$, and keep \$1 in cash.

\step{5}{(d) MM still holds} Total firm value is $50+95=145$ before and $90+55=145$ after, as
\eqref{eq:11V} requires: changing the capital structure did not change the value of the
claims taken together, and investors can undo it (Step~4). The share price rose only because
equal-seniority debt diluted the old creditors, a transfer of \$5 from them to the
shareholders, not value creation.

\intuition MM is a statement about the size of the pie, not about how it is cut. Issuing new
debt that ranks equally with old debt cuts the old creditors' slice; without protective
covenants, shareholders pocket that slice.
\refbib{Grinblatt \& Titman Ch.~14, §14.1 ``The Modigliani-Miller Theorem'' p.~501, §14.2 ``How an Individual Investor Can `Undo' a Firm's Capital Structure Choice'' p.~505, §14.3 ``How Risky Debt Affects the Modigliani-Miller Theorem'' pp.~506--508.}

% =============================================================================
\section{Leveraged recapitalisation with corporate taxes --- Berk 16.18, p.~493}

\notation
\begin{itemize}
  \item $N_0=20$ million shares at $P_0=7.50$; all equity before the announcement.
  \item $D=50$ million of permanent interest-only debt; $\tau_c=0.40$; $P_1$: share price after the debt is issued; $n_R$: shares repurchased (millions).
\end{itemize}

\step{1}{(a) Existing assets} With no debt, assets equal equity: $N_0P_0=(20)(7.50)=150$ million.

\step{2}{(b) Assets after the issue, before the repurchase} The firm holds the old assets, the
\$50 million cash, and the tax shield of permanent debt, $\tau_c D=(0.40)(50)=20$:
\begin{align}
  A_1 &= (150)+(50)+(20) = 220 \text{ million}. \label{eq:12A}
\end{align}

\step{3}{(c) Price and shares repurchased} Equity is $A_1-D=(220)-(50)=170$ million, so
\begin{align}
  P_1 &= \frac{170}{20} = 8.50, \label{eq:12P}\\
  n_R &= \frac{D}{P_1} = \frac{(50)}{(8.50)} = 5.882353 \text{ million shares}. \label{eq:12n}
\end{align}

\step{4}{(d) After the repurchase} The cash is gone:
\begin{center}
\begin{tabular}{lr@{\qquad}lr}
\toprule
Assets & (\$m) & Liabilities and equity & (\$m)\\
\midrule
Existing assets & 150 & Debt & 50\\
Tax shield & 20 & Equity & 120\\
\midrule
Total & 170 & Total & 170\\
\bottomrule
\end{tabular}
\end{center}
With $20-5.882353=14.117647$ million shares left, the price is $120/14.117647=8.50$, unchanged by
the repurchase itself.

\intuition The whole tax benefit, \$20 million or \$1 per share, is capitalised into the price at
the announcement. Shareholders who sell in the buyback and those who stay share it equally, since
the repurchase is at the fair price \$8.50.
\refbib{Berk Ch.~16, §16.3 ``Debt and Taxes'', pp.~473--479 (recapitalising to capture the tax shield).}

% =============================================================================
\section{Derive the MM adjusted cost of capital --- Ross 16.19, p.~578}

\notation
\begin{itemize}
  \item $\mathit{EBIT}$: constant perpetual operating income, fully paid out; $T_C\in[0,1)$: corporate tax rate (no personal taxes).
  \item $D$: perpetual riskless debt with rate $R_D$; $E$: levered equity value; $V=E+D=V_L$: levered firm value; $V_U$: unlevered value.
  \item $R_U$: unlevered cost of capital; $R_E$: cost of levered equity; $\text{CMPC}=\tfrac{E}{V}R_E+\tfrac{D}{V}R_D(1-T_C)$ (Ross's WACC). Assume $E>0$, $V>0$.
\end{itemize}

\textbf{Claim.} $\text{CMPC}=R_U\bigl[1-T_C(D/V)\bigr]$.

\step{1}{Valuation identities (Ross §16.4)} The unlevered firm's perpetual cash flow is
$\mathit{EBIT}(1-T_C)$, so by definition of $R_U$
\begin{align}
  \mathit{EBIT}(1-T_C) = R_U V_U. \label{eq:13VU}
\end{align}
The shield $T_C R_D D$ is a perpetuity as risky as the debt, worth $T_C D$ (cf.\ \eqref{eq:3TS}), so MM I with taxes is
\begin{align}
  V = V_U + T_C D. \label{eq:13MM1}
\end{align}

\step{2}{Returns of the claims} Shareholders receive $(\mathit{EBIT}-R_D D)(1-T_C)$ forever and
debtholders $R_D D$ (a perpetuity worth $D$ at rate $R_D$); valuing the equity stream as a perpetuity gives
\begin{align}
  E\,R_E &= (\mathit{EBIT}-R_D D)(1-T_C). \label{eq:13RE}
\end{align}

\step{3}{Multiply the CMPC by $V$} From the definition,
\begin{align}
  V\cdot\text{CMPC} &= E\,R_E + D\,R_D(1-T_C) \nonumber\\
     &= \bigl[(\mathit{EBIT}-R_D D)(1-T_C)\bigr] + R_D D(1-T_C) && \text{(substituting \eqref{eq:13RE})} \nonumber\\
     &= \mathit{EBIT}(1-T_C) \nonumber\\
     &= \bigl[R_U V_U\bigr] && \text{(substituting \eqref{eq:13VU})} \nonumber\\
     &= R_U\bigl[V - T_C D\bigr] && \text{(substituting \eqref{eq:13MM1})}. \label{eq:13VW}
\end{align}

\step{4}{Divide by $V>0$}
\begin{align}
  \text{CMPC} = R_U\Bigl[1-T_C\frac{D}{V}\Bigr]. \qquad\blacksquare \label{eq:13res}
\end{align}

\step{5}{Same result through MM II} Substituting $\mathit{EBIT}(1-T_C)=R_U(E+D-T_CD)$ (from
\eqref{eq:13VU} and \eqref{eq:13MM1}) into \eqref{eq:13RE} and dividing by $E$ gives MM II with
taxes, $R_E=R_U+(R_U-R_D)\tfrac{D}{E}(1-T_C)$. Then
\begin{align}
  \text{CMPC} &= \tfrac{E}{V}\bigl[R_U+(R_U-R_D)\tfrac{D}{E}(1-T_C)\bigr]+\tfrac{D}{V}R_D(1-T_C) \nonumber\\
     &= \tfrac{E}{V}R_U+\tfrac{D}{V}R_U(1-T_C) \nonumber\\
     &= R_U\Bigl[\tfrac{E+D}{V}-T_C\tfrac{D}{V}\Bigr] = R_U\Bigl[1-T_C\tfrac{D}{V}\Bigr], \nonumber
\end{align}
since $E+D=V$.

\step{6}{Numerical check with Exercise 3} There $T_CD/V=21{,}000/419{,}666.67=0.05003971$, so
$\text{CMPC}=(0.15)(1-0.05003971)=0.14249404$, and $(0.14249404)(419{,}666.67)=59{,}800.00=\mathit{EBIT}(1-T_C)$,
as \eqref{eq:13VW} requires. Directly: $R_E=56{,}290/359{,}666.67=0.156506$ and
$\tfrac{359{,}666.67}{419{,}666.67}(0.156506)+\tfrac{60{,}000}{419{,}666.67}(0.09)(0.65)=0.142494$.

\intuition The levered firm's investors, taken together, receive the unlevered cash flow
$\mathit{EBIT}(1-T_C)$ \emph{plus} the tax saving; the CMPC nets the saving out of the debt cost,
so it must discount the unlevered cash flow to the levered value. Since $V$ exceeds $V_U$ by
$T_CD$, the rate falls below $R_U$ by exactly $R_U T_C D/V$. The formula relies on fixed
perpetual debt; under a constant leverage ratio the shield is riskier and the formula changes.
\refbib{Ross Cap.~16, §16.4, pp.~551--556 (Proposition I with taxes; taxes, WACC and Proposition II).}

% =============================================================================
\section{Miller's gain from leverage with personal taxes --- Amaro de Matos 2.2.7 \#1, p.~58}

\notation
\begin{itemize}
  \item $X$: random operating income per period of two otherwise identical firms, $u$ (unlevered) and $l$ (levered), forever.
  \item $B$: market value of $l$'s perpetual riskless debt, which pays $rB$ per period; $r>0$ is the bond rate.
  \item $\tau_c,\tau_s,\tau_b\in[0,1)$: corporate tax rate (interest deductible), personal tax rate on equity income, personal tax rate on interest income.
  \item $S_u=V_u$, $S_l$: equity values; $V_l=S_l+B$; $\alpha\in(0,1]$: fraction held; $L$: personal loan; $G_L=V_l-V_u$: gain from leverage.
\end{itemize}

\assumption As in Amaro's §2.2 and Appendix answer, individuals borrow and lend at the bond rate
$r$, and personal interest is taxed (lending) or deductible (borrowing) at $\tau_b$; there are no
transaction costs and no arbitrage.

\step{1}{Payoff of the levered equity} The firm pays corporate tax on $X-rB$, and the holder pays $\tau_s$ on what she receives:
\begin{align}
  Y_l = \alpha(X-rB)(1-\tau_c)(1-\tau_s). \label{eq:14Yl}
\end{align}

\step{2}{Payoff of the unlevered equity}
\begin{align}
  Y_u = \alpha X(1-\tau_c)(1-\tau_s). \label{eq:14Yu}
\end{align}

\step{3}{Size of the personal loan} A loan $L$ at $r$ costs $rL(1-\tau_b)$ per period after tax.
Choose $L$ so this equals the debt-service part of \eqref{eq:14Yl}:
\begin{align}
  rL(1-\tau_b) &= \alpha rB(1-\tau_c)(1-\tau_s) \nonumber\\
  L &= \alpha B\,\frac{(1-\tau_c)(1-\tau_s)}{1-\tau_b}. \label{eq:14L}
\end{align}

\step{4}{Replication} Buying $\alpha$ of $u$ and borrowing $L$ pays, by \eqref{eq:14Yu} and \eqref{eq:14L},
\begin{align}
  Y_u - rL(1-\tau_b) &= \alpha X(1-\tau_c)(1-\tau_s)-\bigl[\alpha rB(1-\tau_c)(1-\tau_s)\bigr] \nonumber\\
     &= \alpha(X-rB)(1-\tau_c)(1-\tau_s) = Y_l, \nonumber
\end{align}
in every state and every period. Its cost today is $\alpha S_u-L$.

\step{5}{Law of one price} If $\alpha S_l>\alpha S_u-L$, a holder of $l$'s shares sells them, builds the
replica and keeps the difference with unchanged future income. If $\alpha S_l<\alpha S_u-L$, a holder
of $u$'s shares sells them, buys $\alpha$ of $l$'s shares and lends $L$ (after-tax interest
$rL(1-\tau_b)$, restoring $Y_u$ by Step~4 read backwards), again pocketing the difference.
No arbitrage therefore requires
\begin{align}
  \alpha S_l &= \alpha S_u - L \nonumber\\
  S_l &= S_u - B\,\frac{(1-\tau_c)(1-\tau_s)}{1-\tau_b} && \text{(substituting \eqref{eq:14L}, dividing by } \alpha). \label{eq:14Sl}
\end{align}

\step{6}{Firm values and the gain} Adding $B$ to both sides of \eqref{eq:14Sl} and using $V_l=S_l+B$, $V_u=S_u$:
\begin{align}
  V_l &= V_u + B - B\,\frac{(1-\tau_c)(1-\tau_s)}{1-\tau_b}, \nonumber\\
  G_L = V_l - V_u &= \Bigl[1-\frac{(1-\tau_c)(1-\tau_s)}{1-\tau_b}\Bigr]B. \qquad\blacksquare \label{eq:14GL}
\end{align}
The gain goes to the stockholders who decide to lever: they collect $B$ at issue and keep shares
worth $S_l$, a total of $V_l$ instead of $V_u$.

\step{7}{Special cases} (i) $\tau_s=\tau_b$ (including no personal taxes): $G_L=\tau_cB$, the MM
(1963) result, Amaro (2.7). (ii) Miller's equilibrium (2.11), $(1-\tau_c)(1-\tau_s)=1-\tau_b$:
$G_L=0$, so no single firm gains from leverage. (iii) $(1-\tau_c)(1-\tau_s)>1-\tau_b$: $G_L<0$,
and the firm should not borrow.

\bookans The final expression matches the book's answer (p.~255). The intermediate line printed
there, ``$\alpha S_l=\alpha S_u+\alpha B\frac{(1-\tau_c)(1-\tau_s)}{(1-\tau_b)}$'', has a sign typo:
it must be a minus, as in \eqref{eq:14Sl}, since the investor \emph{borrows} that amount and the
book's own next line, $V_l=V_u+B[1-\cdots]$, follows only with the minus.

\intuition Debt moves a dollar of income from the corporate-plus-equity tax channel, where it
keeps $(1-\tau_c)(1-\tau_s)$, to the interest channel, where it keeps $(1-\tau_b)$. Leverage
creates value only to the extent that the second channel is less taxed than the first.
\refbib{Amaro de Matos Ch.~2, §2.2 ``Personal and Corporate Taxation'', pp.~48--52 (Miller equilibrium, $G_L$ on p.~52); Appendix C.2.2, p.~255.}

% =============================================================================
\section{MM I with perpetual debt, and when borrowing and lending rates differ --- Amaro de Matos 2.1.5 \#3--\#4, p.~48}

\notation
\begin{itemize}
  \item $X$: random operating income per period (the time average of the stream), identical for firms $u$ (unlevered) and $l$ (levered); any distribution.
  \item $B$: market value of $l$'s perpetual riskless debt, paying $rB$ per period; $r$: corporate interest rate.
  \item $\tau\in[0,1)$: corporate tax rate, interest deductible ($\tau=0$ is the no-tax case); no personal taxes.
  \item $V_u=S_u$, $V_l=S_l+B$; $\alpha\in(0,1]$: fraction of a firm held by an investor.
  \item \#4 only: $\bar r$ and $\underline r$: individual borrowing and lending rates, $\bar r>r>\underline r>0$.
\end{itemize}

\assumption In \#3 individuals borrow and lend at $r$, as firms do; there are no transaction
costs. In \#4, as in the book's answer, $\tau=0$ and individuals can lend only at $\underline r$
(they cannot buy $l$'s bonds at yield $r$).

\subsection*{\#3 --- MM I with and without taxes}

\step{1}{Payoffs} Shareholders of $u$ and $l$ receive, each period,
\begin{align}
  Y_u = X(1-\tau), \qquad Y_l = (X-rB)(1-\tau). \label{eq:15Y}
\end{align}

\step{2}{Replicating $l$'s equity} Buy $\alpha$ of $u$ and borrow $\alpha(1-\tau)B$ at $r$. By \eqref{eq:15Y} the income is
\begin{align}
  \alpha X(1-\tau) - r\,\alpha(1-\tau)B = \alpha(X-rB)(1-\tau) = \alpha Y_l, \nonumber
\end{align}
at a cost of $\alpha S_u-\alpha(1-\tau)B$ today.

\step{3}{Upper bound} If $S_l>S_u-(1-\tau)B$, a holder of $\alpha$ of $l$'s shares sells them for
$\alpha S_l$, builds the replica of Step~2 and keeps, with identical future income,
\begin{align}
  \alpha\bigl[S_l-S_u+(1-\tau)B\bigr] &= \alpha\bigl[(V_l-B)-V_u+(1-\tau)B\bigr] \nonumber\\
      &= \alpha\bigl[V_l-V_u-\tau B\bigr] > 0, \nonumber
\end{align}
an arbitrage. Hence $V_l\le V_u+\tau B$.

\step{4}{Replicating $u$'s equity} Buy $\alpha$ of $l$'s shares and $\alpha(1-\tau)$ of its debt (or lend $\alpha(1-\tau)B$ at $r$):
\begin{align}
  \alpha(X-rB)(1-\tau)+\alpha(1-\tau)rB = \alpha X(1-\tau) = \alpha Y_u, \nonumber
\end{align}
at a cost of $\alpha\bigl[S_l+(1-\tau)B\bigr]$.

\step{5}{Lower bound} If $S_u>S_l+(1-\tau)B$, i.e.\ $V_l<V_u+\tau B$, a holder of $u$'s shares sells them,
buys the replica of Step~4 and keeps $\alpha[V_u+\tau B-V_l]>0$. Hence $V_l\ge V_u+\tau B$.

\step{6}{MM I} Steps 3 and 5 give
\begin{align}
  V_l = V_u + \tau B \quad\text{(with taxes)}, \qquad V_l = V_u \quad (\tau=0). \qquad\blacksquare \label{eq:15MM1}
\end{align}

\bookans Matches the book's answer (pp.~252--253), $V_u+\tau B=V_l$. The book's switching
portfolio invests all the proceeds in the other firm instead of replicating exactly, which leads
to the same two inequalities.

\subsection*{\#4 --- Different borrowing and lending rates ($\tau=0$)}

\step{7}{Replicating $l$'s equity with costly borrowing} The replica must still hold $\alpha$ of $u$
(the only other claim on $X$), so its debt service has to equal $\alpha rB$ at the rate $\bar r$:
\begin{align}
  \bar r\,\bar L = \alpha rB \quad\Longrightarrow\quad \bar L = \alpha B\,\frac{r}{\bar r}, \label{eq:15Lbar}
\end{align}
which gives income $\alpha X-\alpha rB=\alpha Y_l$ at a cost of $\alpha\bigl(S_u-B\,r/\bar r\bigr)$.
(Borrowing the full $\alpha B$ would cost $\alpha\bar rB>\alpha rB$ per period: not a replica.)

\step{8}{Upper bound} No arbitrage against an overvalued $l$ requires $S_l\le S_u-B\,r/\bar r$; adding $B$,
\begin{align}
  V_l \le V_u + B\Bigl(1-\frac{r}{\bar r}\Bigr) = V_u + B\,\frac{\bar r-r}{\bar r}. \label{eq:15up}
\end{align}

\step{9}{Replicating $u$'s equity with poor lending} Buy $\alpha$ of $l$'s shares and lend $\underline L$ at $\underline r$ with
\begin{align}
  \underline r\,\underline L = \alpha rB \quad\Longrightarrow\quad \underline L = \alpha B\,\frac{r}{\underline r}>\alpha B, \label{eq:15Llow}
\end{align}
giving income $\alpha(X-rB)+\alpha rB=\alpha Y_u$ at a cost of $\alpha\bigl(S_l+B\,r/\underline r\bigr)$.

\step{10}{Lower bound} No arbitrage against an overvalued $u$ requires $S_u\le S_l+B\,r/\underline r$, i.e.\ $V_u\le V_l-B+B\,r/\underline r$:
\begin{align}
  V_l \ge V_u - B\,\frac{r-\underline r}{\underline r}. \label{eq:15low}
\end{align}

\step{11}{What happens to MM1} Combining \eqref{eq:15up} and \eqref{eq:15low},
\begin{align}
  V_u - B\,\frac{r-\underline r}{\underline r} \;\le\; V_l \;\le\; V_u + B\,\frac{\bar r-r}{\bar r}. \qquad\blacksquare \label{eq:15band}
\end{align}
MM1 no longer pins down $V_l$: arbitrage only confines it to a band around $V_u$ whose width grows
with the debt $B$ and with the spreads. With $\bar r=r+\epsilon$, $\underline r=r-\epsilon$ the band is
$V_u-\epsilon B/(r-\epsilon)\le V_l\le V_u+\epsilon B/(r+\epsilon)$, and it collapses to $V_l=V_u$
as $\epsilon\to0$. Leverage can now raise firm value (up to the upper bound): the firm borrows
more cheaply than individuals, so corporate leverage is a service investors cannot fully copy.
If individuals could buy $l$'s bonds at yield $r$, Step~4 would apply and the lower bound would
tighten to $V_l\ge V_u$.

\bookans Differs from the book's printed band (p.~254), $V_u(1-\epsilon B_l/X)\le V_l\le V_u(1+\epsilon B_l/X)$.
The qualitative conclusion is the same (a band around $V_u$, proportional to $\epsilon B$, vanishing at
$\epsilon=0$), but the bounds differ. The book's portfolios keep the investor's capital fixed (all
proceeds go into the other firm, and the full principal $B_l$ is borrowed or lent), so the two
incomes differ by a multiple of the random $X$ minus a fixed amount. Its inequality therefore holds
only for a given value of $X$ (or in expectation), not state by state. The portfolios above match
incomes exactly, so \eqref{eq:15band} is a true no-arbitrage band for any distribution of $X$. The
two coincide only if $X/V_u$ happens to equal $\bar r$ (upper bound) or $\underline r$ (lower bound).

\intuition MM1 rests on investors being able to copy corporate leverage on the same terms. When
personal borrowing is dearer and lending cheaper than corporate debt, homemade leverage is an
imperfect substitute: arbitrage still limits mispricing, but only to within the cost of the
personal-borrowing wedge.
\refbib{Amaro de Matos Ch.~2, §2.1.1--2.1.3, pp.~40--46 (MM1, MM2, MM with taxes); Appendix C.2.1, pp.~252--254.}

\end{document}
